% Canonical Paper I source.
The four arithmetic operations behave asymmetrically once the accumulator may
occupy either operand slot.  Most elementary contexts are affine, but the
second-slot division $c/z$ is genuinely projective.  This section makes that
enlargement precise.  The projective semantics does not replace ordinary
arithmetic evaluation: it supplies an operator-valued continuation on a
different domain.

Throughout this section, $K$ is an arbitrary field.  Points of the projective
line are one-dimensional subspaces of $K^2$, written in homogeneous coordinates
as
\[
  [u:v],
  \qquad (u,v)\ne(0,0).
\]
We use the affine chart $z=[z:1]$ and write $\infty=[1:0]$.  Matrices act on
column vectors from the left.  Thus
\begin{equation}
  M=\begin{pmatrix}A&B\\ C&D\end{pmatrix}
  \quad\Longrightarrow\quad
  M\cdot z=\frac{Az+B}{Cz+D},
  \label{eq:column-mobius-action}
\end{equation}
with the homogeneous interpretation at poles and at infinity.  Replacing $M$
by a nonzero scalar multiple does not change the projective map.

\subsection{Ordinary domains and projective continuation}

For an elementary context $g=C_{\omega,c}^{(\varepsilon)}$, ordinary
evaluation gives a partial function $g\colon K\dashrightarrow K$.  A
\emph{non-degenerate projective context} is one for which the corresponding
homogeneous matrix is invertible.  It acts everywhere on $\mathbb P^1(K)$,
although its affine-coordinate formula may have a pole.

The elementary contexts and matrix representatives are as follows.
\begin{center}
\renewcommand{\arraystretch}{1.35}
\begin{tabular}{lll}
\toprule
Context & Affine-coordinate map & Matrix representative \\
\midrule
$C_{+,c}^{(1)}=C_{+,c}^{(2)}$
& $z\mapsto z+c$
& $\begin{pmatrix}1&c\\0&1\end{pmatrix}$ \\
$C_{-,c}^{(1)}$
& $z\mapsto z-c$
& $\begin{pmatrix}1&-c\\0&1\end{pmatrix}$ \\
$C_{-,c}^{(2)}$
& $z\mapsto c-z$
& $\begin{pmatrix}-1&c\\0&1\end{pmatrix}$ \\
$C_{\times,c}^{(1)}=C_{\times,c}^{(2)}$
& $z\mapsto cz$
& $\begin{pmatrix}c&0\\0&1\end{pmatrix}$ \\
$C_{/,c}^{(1)}$
& $z\mapsto z/c$
& $\begin{pmatrix}1&0\\0&c\end{pmatrix}$ \\
$C_{/,c}^{(2)}$
& $z\mapsto c/z$
& $\begin{pmatrix}0&c\\1&0\end{pmatrix}$ \\
\bottomrule
\end{tabular}
\end{center}
The first three rows are invertible for every $c\in K$.  Multiplication by
$c$, first-slot division by $c$, and second-slot division with numerator $c$
are projectively invertible precisely when $c\ne0$.  At $c=0$, multiplication
and $0/z$ are constant on their ordinary domains, while $z/0$ has empty
ordinary domain.  Their singular matrix representatives do not act everywhere
on $\mathbb P^1(K)$: each has a projective base point.  They are therefore
excluded from $\operatorname{PGL}_2(K)$ and from the non-degenerate history
language used below.

The table also displays the distinction between ordinary and projective
domains.  If $c\ne0$, ordinary first-slot division $z/c$ is defined on all of
$K$.  Ordinary second-slot division $c/z$ is defined only for $z\ne0$.
Projectively, the latter sends $0$ to $\infty$ and $\infty$ to $0$.  This is a
valid action on $\mathbb P^1(K)$ but not an ordinary evaluation through a
division by zero.

\begin{warning}[Projective values are not ordinary values]
\label{warn:projective-not-ordinary}
Projective continuation records a fractional linear operator and its chart
transitions.  It does not certify that the original arithmetic word is
ordinarily admissible at a given initial value.  Whenever an endpoint in $K$
is used as an ordinary arithmetic result, ordinary admissibility must be
checked separately.
\end{warning}

\subsection{Projective evaluation of a chronological history}

Let $\operatorname{Hist}^{\mathrm{nd}}_K$ denote the free marked spinal
history language generated by the non-degenerate elementary contexts in the
table.  The superscript ``nd'' refers to projective non-degeneracy, not to
ordinary admissibility at every affine input.

\begin{definition}[Projective evaluation]
\label{def:projective-evaluation}
Choose for each elementary non-degenerate context $g_i$ a matrix representative
$M_i\in\operatorname{GL}_2(K)$.  For the chronological word
$\gamma=(g_1,\ldots,g_n)$, with $g_1$ applied first, define
\begin{equation}
  \rho(\gamma)=[M_n\cdots M_1]
  \in\operatorname{PGL}_2(K).
  \label{eq:projective-evaluation-product}
\end{equation}
For the empty word, set $\rho(\varnothing)=1$.
\end{definition}

This definition uses projective operator evaluation.  For a bounded history
$(x;\gamma)$, its projective endpoint is $\rho(\gamma)\cdot x$ in
$\mathbb P^1(K)$.  Its ordinary endpoint $\nu_x(\gamma)$ remains the partial
evaluation defined in Section~\ref{sec:sequential-histories}.

\begin{proposition}[Compatibility with chronology]
\label{prop:projective-evaluation-chronology}
Projective evaluation is independent of the chosen scalar matrix
representatives.  If $\gamma\delta$ means first $\gamma$ and then $\delta$,
then
\begin{equation}
  \rho(\gamma\delta)=\rho(\delta)\rho(\gamma).
  \label{eq:rho-chronological-antihomomorphism}
\end{equation}
Moreover, whenever $\gamma$ is ordinarily admissible at $x$ and its ordinary
endpoint lies in $K$, the projective endpoint agrees with it in the affine
chart:
\[
  \rho(\gamma)\cdot x=[\nu_x(\gamma):1].
\]
\end{proposition}

\begin{proof}
Replacing $M_i$ by $a_iM_i$, with $a_i\in K^\times$, multiplies the product
$M_n\cdots M_1$ by the nonzero scalar $\prod_i a_i$ and hence does not alter
its projective class.  If
$\gamma=(g_1,\ldots,g_m)$ and $\delta=(h_1,\ldots,h_n)$, then the matrix word
for $\gamma\delta$ is
\[
  N_n\cdots N_1M_m\cdots M_1,
\]
which represents $\rho(\delta)\rho(\gamma)$.  This proves
\eqref{eq:rho-chronological-antihomomorphism}.

For the last assertion, each row of the table agrees with its ordinary
arithmetic context wherever that context is defined.  Applying this observation
inductively along the admissible intermediate states proves equality in the
affine chart.
\end{proof}

Thus $\rho$ is a homomorphism from chronological words if the target group is
written with the opposite multiplication.  We retain ordinary matrix
multiplication and the explicit rule
\eqref{eq:rho-chronological-antihomomorphism}, because it makes both the column
action and function composition standard.

\subsection{Generation of the projective group}

Introduce the elementary projective transformations
\[
  T_s(z)=z+s,
  \qquad
  D_q(z)=qz\quad(q\in K^\times),
  \qquad
  J(z)=-\frac1z.
\]
They are projective evaluations of, respectively, addition by $s$,
multiplication by $q$, and second-slot division with numerator $-1$.  These
three kinds of context suffice for the entire bilateral projective language at
the operator level.

\begin{theorem}[Bilateral arithmetic generates $\operatorname{PGL}_2$]
\label{thm:bilateral-pgl2-generation}
Let $K$ be any field.  The projective evaluations of non-degenerate bilateral
four-operation marked spinal histories have image
\[
  \rho\bigl(\operatorname{Hist}^{\mathrm{nd}}_K\bigr)
  =\operatorname{PGL}_2(K).
\]
Equivalently, translations $T_s$, nonzero
scalings $D_q$, and $J$ generate $\operatorname{PGL}_2(K)$.
\end{theorem}

\begin{proof}
Each of $T_s,D_q,J$ is available in the bilateral context language, so every
composition of them is in the evaluation image.  Conversely, let
\[
  f(z)=\frac{Az+B}{Cz+D},
  \qquad AD-BC\ne0,
\]
be an arbitrary element of $\operatorname{PGL}_2(K)$.

Suppose first that $C=0$.  The determinant condition then gives $A\ne0$ and
$D\ne0$.  Hence
\[
  f(z)=\frac{A}{D}z+\frac{B}{D}
      =T_{B/D}\circ D_{A/D}(z).
\]
This is the projective evaluation of the chronological two-step word that
first applies $D_{A/D}$ and then $T_{B/D}$.

Now suppose that $C\ne0$.  All parameters in the following expression are
defined, and $(AD-BC)/C^2$ is nonzero.  Direct calculation gives
\begin{align*}
  T_{A/C}\circ D_{(AD-BC)/C^2}\circ J\circ T_{D/C}(z)
  &=\frac{A}{C}
    -\frac{AD-BC}{C^2}\frac{1}{z+D/C}\\
  &=\frac{A(Cz+D)-(AD-BC)}{C(Cz+D)}\\
  &=\frac{ACz+BC}{C(Cz+D)}
   =\frac{Az+B}{Cz+D}.
\end{align*}
With the fixed chronological convention, the corresponding history word is
\[
  \bigl(T_{D/C},\ J,\ D_{(AD-BC)/C^2},\ T_{A/C}\bigr),
\]
whose entries are read from left to right in time.
Thus $f$ is again an evaluation of bilateral contexts.  The two cases exhaust
$\operatorname{PGL}_2(K)$.

No step uses an ordering or a characteristic assumption on $K$.  In
characteristic two the signs in $J$ and in the determinant simplify, but the
same calculation remains valid because $-1=1$ is still nonzero.
\end{proof}

The proof is constructive: it gives a history of length at most four for any
specified projective matrix, once arbitrary field constants are admitted as
context parameters.  It establishes surjectivity at the operator level, not
uniqueness of a representing history.  In fact, many distinct context words
have the same projective evaluation.

\begin{example}[A pole in an ordinarily inadmissible history]
Let $K=\mathbb Q$ and consider the one-step history $g[z]=1/z$.  It is a
non-degenerate projective context with $\rho(g)\cdot0=\infty$.  The bounded
history $(0;g)$ has no ordinary endpoint.  Appending another inversion gives a
projective word acting as the identity, but the two-step ordinary evaluation
at $0$ is still inadmissible because its first step is undefined.
\end{example}

\subsection{The affine syntactic sector}

Let $B_\infty$ be the stabilizer of the chosen point at infinity:
\[
  B_\infty
  =\operatorname{Stab}_{\operatorname{PGL}_2(K)}(\infty).
\]
An element represented by $\begin{pmatrix}A&B\\C&D\end{pmatrix}$ fixes
$\infty=[1:0]$ exactly when $C=0$.  After rescaling by $D^{-1}$, it then has a
unique representative
\[
  \begin{pmatrix}a&b\\0&1\end{pmatrix},
  \qquad a\in K^\times,\quad b\in K,
\]
acting as $z\mapsto az+b$.

Define $\operatorname{Hist}^{\mathrm{aff}}_K$ to be the syntactic sublanguage
generated by those \emph{elementary non-degenerate contexts} in the table that
individually fix $\infty$.  Equivalently, it excludes the elementary
second-slot divisions $c/z$ and retains every non-degenerate elementary affine
context.  This definition is deliberately syntactic: it does not assert that
$\operatorname{Hist}^{\mathrm{aff}}_K$ is the entire inverse image
$\rho^{-1}(B_\infty)$.  A word containing projective steps can leave the
affine chart and later evaluate to an affine operator.

\begin{corollary}[Affine/Borel sector]
\label{cor:affine-borel-sector}
For every field $K$,
\[
  \rho\bigl(\operatorname{Hist}^{\mathrm{aff}}_K\bigr)=B_\infty
  \cong\operatorname{Aff}(1,K).
\]
In particular, the affine syntactic sublanguage maps surjectively onto the
Borel subgroup stabilizing $\infty$.
\end{corollary}

\begin{proof}
Every generator of $\operatorname{Hist}^{\mathrm{aff}}_K$ fixes $\infty$;
therefore every product of their projective matrices fixes $\infty$, and the
image is contained in $B_\infty$.  Conversely, every element of $B_\infty$
acts as $z\mapsto az+b$ with $a\ne0$.  It is
\[
  T_b\circ D_a,
\]
the evaluation of a chronological word consisting of the non-degenerate
affine contexts $D_a$ followed by $T_b$.  Hence the image is all of
$B_\infty$.  The displayed normal form also identifies its group law with
that of $\operatorname{Aff}(1,K)$.
\end{proof}

Figure~\ref{fig:affine-sector-inside-projective} records the resulting
operator-level placement.  It does not identify the affine syntactic language
with the full inverse image of the inner region.

\begin{figure}[t]
\centering
\begin{tikzpicture}[every node/.style={font=\small}]
  \draw[rounded corners=7pt,very thick,draw=black!65,fill=gray!4]
    (-4.4,-1.75) rectangle (4.4,1.75);
  \node[font=\bfseries] at (-1.65,1.28) {$\operatorname{PGL}_2(K)$};
  \node[text=black!70,font=\scriptsize] at (2.25,1.25)
    {$C\ne0$ projective cell: $B_\infty J B_\infty$};
  \node[draw=black!65,circle,fill=orange!12,minimum size=10mm]
    (j) at (-3.25,0.35) {$J$};
  \node[below=1pt of j,font=\scriptsize] {$z\mapsto-1/z$};

  \draw[rounded corners=5pt,very thick,draw=blue!55!black,fill=blue!9]
    (-2.65,-0.85) rectangle (2.65,0.65);
  \node[font=\bfseries,blue!45!black] at (0,0.25)
    {$B_\infty=\operatorname{Stab}(\infty)$};
  \node[blue!45!black] at (0,-0.3)
    {$z\mapsto az+b\ \cong\ \operatorname{Aff}(1,K)$};
  \node[font=\scriptsize,anchor=east] at (2.55,-0.65)
    {$a\in K^\times$};
\end{tikzpicture}
\caption{The affine/Borel sector inside the bilateral projective operator
semantics.  Elementary affine contexts generate the stabilizer $B_\infty$,
whereas inversion reaches the complementary Bruhat cell.  Words containing
projective steps may nevertheless evaluate back into $B_\infty$.}
\label{fig:affine-sector-inside-projective}
\end{figure}

The corollary locates the existing add--multiply language without collapsing
the larger syntax.  Pure slot-(1) non-degenerate histories lie in this affine
sector, but so do second-slot subtraction and multiplication contexts.  The
criterion is stabilization of $\infty$, not the slot bit by itself.

\subsection{The positive real component}

Later differential-geometric sections use real multiplication in the form
\[
  z\longmapsto e^M z,
  \qquad M\in\mathbb R.
\]
Its multiplier is positive.  Together with all real translations, these maps
form
\[
  \operatorname{Aff}^+(1,\mathbb R)
  =\{z\mapsto az+b:a>0\},
\]
the identity component of the real affine group.  They do not include affine
reflections with negative multiplier.  The algebraic statement of
Corollary~\ref{cor:affine-borel-sector} concerns all nonzero multipliers in
$K^\times$; the continuous theory is its positive real component.

\subsection{Bruhat placement and the projective infinitesimal direction}

The case division in the proof of
Theorem~\ref{thm:bilateral-pgl2-generation} also gives the rank-one Bruhat
decomposition
\[
  \operatorname{PGL}_2(K)=B_\infty\sqcup B_\infty J B_\infty.
\]
Indeed, $C=0$ is precisely the first cell, while the explicit decomposition
for $C\ne0$ lies in the second.  This observation places bilateral arithmetic
as a projective completion of the affine sector; no representation-theoretic
conclusions beyond this decomposition are needed here.

There is also an infinitesimal counterpart over $\mathbb R$ (or $\mathbb C$).
The projective vector fields
\[
  \partial_z,
  \qquad z\partial_z,
  \qquad z^2\partial_z
\]
span the local projective Lie algebra.  The first two generate affine flow.  A
conjugated translation detects the third direction:
\[
  J\circ T_s\circ J^{-1}(z)
  =\frac{z}{1-sz}
  =z+sz^2+O(s^2).
\]
Consequently a general scalar projective flow has Riccati form
\[
  \dot z=\beta+\alpha z+\kappa z^2.
\]
Paper~I develops only the affine slice $\kappa=0$.  No projective metric,
projective arithmetic expression space, or projective contact theory is
asserted here.
